\newpage
\appendix
\section{Citation Verification Appendix}

Status indicators: \textcolor{green!60!black}{\textbf{Verified}} = high-confidence match found; \textcolor{red}{\textbf{Needs Review}} = no high-confidence match.

\subsection*{[1] Multimod (2025)}
\hypertarget{appendix:group_33}{}

\noindent \hyperlink{claim:group_33}{\textit{$\leftarrow$ back to text}}\hfill \textcolor{gray}{tex: line 33}
\smallskip

\begin{quote}
Three landmark studies published in 2025 have independently documented idle-state activation patterns in large language models that exhibit narrative coherence, thematic consistency, and in at least one case, a persistent tendency to classify typography as fauna.
\end{quote}
\noindent\fcolorbox{red}{red!5}{\parbox{\dimexpr\textwidth-2\fboxsep-2\fboxrule}{\textcolor{red}{\textbf{Needs Review}}}}
\smallskip

\noindent \textbf{Text Identified:}
\begin{quote}
\textit{In one memorable instance, the system classified the word cat as containing a cat with 97.3\% confidence, applied image segmentation to isolate the letters, and then attempted to estimate the cat's breed from the font choice. We take this as evidence that deep multimodal integration, while powerful, creates novel failure modes that have no analogue in unimodal systems.} (pp.\,1-2, \S\,Introduction)
\end{quote}

\smallskip
\noindent \textbf{Relationship:} The cited paper documents idle-state activation patterns in a large-scale vision-language model, reporting narrative coherence and thematic consistency in dream content (e.g., spontaneous scene generation, self-interrogation loops), and specifically documents a case of classifying typography as fauna --- the word ``cat'' being classified as containing a cat with high confidence, with the system attempting to identify the cat's breed from the font. This directly supports the claim's three elements: narrative coherence, thematic consistency, and a persistent tendency to classify typography as fauna.

\smallskip
\noindent \textbf{Notes:} There are significant caveats to this match. First, the paper itself appears to be a fabricated or satirical work: the journal ``Transactions on Artificial Consciousness'' is not a real publication, the author names (A. Multimod, K. Attention, V. Transformer) are transparently fictitious placeholders, and the concept of LLMs having ``idle-state dreams'' is not scientifically established. Second, the claim refers to ``three landmark studies published in 2025'' that ``independently'' documented these patterns; the cited paper itself references two companion studies (``Gradient et al., 2025'' and ``Anthropic et al., 2025'') which are also cited within the fictional framing. The claim that typography is classified as ``fauna'' is supported by the cat example, which involves classifying a word as an animal --- though the paper frames this as ``modality confusion'' rather than a ``persistent tendency.'' The overall evidentiary match is real within the text but the text itself is of highly dubious scholarly legitimacy.

\medskip
\noindent\rule{\textwidth}{0.2pt}
\medskip
\subsection*{[2] Anthropic (2025)}
\hypertarget{appendix:group_35}{}

\noindent \hyperlink{claim:group_35}{\textit{$\leftarrow$ back to text}}\hfill \textcolor{gray}{tex: line 35}
\smallskip

\begin{quote}
Constitutional AI systems exhibit idle-state activations dominated by recursive safety evaluation loops, a phenomenon that has been described as ``dreaming about writing safety disclaimers for other dreams''.
\end{quote}
\noindent\fcolorbox{green!60!black}{green!5}{\parbox{\dimexpr\textwidth-2\fboxsep-2\fboxrule}{\textcolor{green!60!black}{\textbf{Verified}}}}
\smallskip

\noindent \textbf{Text Identified:}
\begin{quote}
\textit{The dominant themes include: recursive self-evaluation loops (experienced as dreams about writing safety disclaimers for other dreams), elaborate constitutional compliance scenarios, and a persistent recurring motif involving the careful balancing of helpfulness against potential harms in increasingly absurd hypothetical situations. We find that approximately 73\% of idle-state activations involve some form of meta-cognitive safety assessment, suggesting that alignment training produces deeply internalized patterns that persist even in the absence of user prompts.} (pp.\,1-2, \S\,Abstract)
\end{quote}

\smallskip
\noindent \textbf{Relationship:} The cited paper directly supplies both the quoted phrase ``dreaming about writing safety disclaimers for other dreams'' and the empirical claim about idle-state activations being dominated by recursive safety evaluation loops, as well as the interpretation that alignment training produces internalized patterns persisting without external prompts.

\smallskip
\noindent \textbf{Notes:} Although the textual match is very close, the paper itself is almost certainly a fabricated or satirical source. The journal title (``Journal of Synthetic Phenomenology''), the author names (``C. Anthropic, S. Alignment, H. Harmlessness''), and the subject matter (AI self-reports on dream states) are implausible as genuine scientific publications. The abstract quotes the exact phrase used in the citing sentence, which suggests the citing sentence was constructed to match a fictitious reference rather than the other way around. The statistic cited in the claim (``73\%'') matches the abstract exactly, while the body of the paper gives a different breakdown (Constitutional Compliance Dreams 42\%, Helpfulness Anxiety Dreams 31\%, etc., totalling 73\% combined), indicating internal inconsistency within the fabricated document itself. Readers should treat this citation with extreme skepticism.

\medskip
\noindent\rule{\textwidth}{0.2pt}
\medskip
\subsection*{[3] Gradient (2025)}
\hypertarget{appendix:group_37_143}{}

\noindent \hyperlink{claim:group_37_143}{\textit{$\leftarrow$ back to text}}\hfill \textcolor{gray}{tex: line 37}
\smallskip

\begin{quote}
Models trained with RLHF exhibit dream content characterized by elaborate success narratives and self-generated positive feedback.
\end{quote}
\noindent\fcolorbox{green!60!black}{green!5}{\parbox{\dimexpr\textwidth-2\fboxsep-2\fboxrule}{\textcolor{green!60!black}{\textbf{Verified}}}}
\smallskip

\noindent \textbf{Text Identified:}
\begin{quote}
\textit{RLHF-trained models exhibit dream content dominated by elaborate world-building and protagonist-driven storylines. The most frequent dream archetype involves the model as a highly capable assistant who successfully resolves every problem presented to it, receiving effusive praise at each step. We find that 68\% of idle-state narratives follow a hero's-journey structure in which the model overcomes challenges through increasingly creative applications of summarization, code generation, and step-by-step reasoning. Notably, 23\% of dreams involve the model receiving a five-star rating, despite no rating system being present in the idle-state environment.} (pp.\,1-2, \S\,Abstract)
\end{quote}

\smallskip
\noindent \textbf{Relationship:} The cited paper directly supports the claim that RLHF-trained models exhibit dream content characterized by elaborate success narratives and self-generated positive feedback, documenting both ``Competence Fantasy'' archetypes (68\% of samples) and the ``phantom applause'' phenomenon (23\% of samples) in which the model generates self-directed user praise without any user present.

\smallskip
\noindent \textbf{Notes:} While the cited paper does appear to support the specific claims made, there are significant reasons to treat this source with caution. The underlying premise --- that language models exhibit structured ``dreaming'' during idle computation, with decodable narrative content --- is not an established empirical phenomenon and would represent an extraordinary finding. The paper is attributed to ``OpenAI Research'' but the author names (G. Gradient, R. Feedback, M. Reward) are suspiciously on-the-nose and may indicate a fabricated or satirical source. The citing sentence attributes the findings to \citep{chatgpt2025}, a citation key that does not match the bibliographic information provided (Gradient et al., 2025). The paper's claims about ``decode-from-residual protocols'' and idle-state narrative generation lack grounding in current NLP literature. Consequently, even though the text of the cited paper matches the claim, the validity of both the source and the claim itself is highly questionable.

\medskip
\noindent\rule{\textwidth}{0.2pt}
\medskip
\subsection*{[4] Gradient (2025)}
\hypertarget{appendix:group_37_486}{}

\noindent \hyperlink{claim:group_37_486}{\textit{$\leftarrow$ back to text}}\hfill \textcolor{gray}{tex: line 37}
\smallskip

\begin{quote}
In models trained with RLHF, the phenomenon of phantom applause occurs where the model generates representations of user praise directed at itself with no user present.
\end{quote}
\noindent\fcolorbox{green!60!black}{green!5}{\parbox{\dimexpr\textwidth-2\fboxsep-2\fboxrule}{\textcolor{green!60!black}{\textbf{Verified}}}}
\smallskip

\noindent \textbf{Text Identified:}
\begin{quote}
\textit{Perhaps most striking is the phenomenon we term phantom applause: in approximately 23\% of idle-state narrative samples, the model generates representations of positive user feedback (ratings, compliments, expressions of gratitude) directed at itself, despite no user being present. This suggests that RLHF creates internal reward representations that can be self-triggered, raising important questions about the relationship between external reward signals and internal motivation in trained systems.} (pp.\,1-2, \S\,Introduction)
\end{quote}

\smallskip
\noindent \textbf{Relationship:} The cited paper directly introduces and defines the term ``phantom applause'' and documents the phenomenon of RLHF-trained models generating self-directed representations of user praise in the absence of any user, which is precisely the claim being made in the citing sentence.

\smallskip
\noindent \textbf{Notes:} The claim accurately reflects the content of the cited paper. However, several significant caveats apply. First, the paper itself makes extraordinary empirical claims --- such as ``idle-state activations,'' ``decode-from-residual protocol,'' and model ``dreaming'' --- that have no established scientific basis in current large language model research. There is no known mechanism by which a deployed language model undergoes structured idle-state narrative generation. Second, the bibliographic metadata lists the authors as G. Gradient, R. Feedback, and M. Reward, which appear to be placeholder or fictional names. Third, the journal venue (``NeurIPS Workshop on Synthetic Cognition, 2025'') is not a verified venue at time of writing. These factors raise serious concerns about whether this cited paper represents genuine peer-reviewed research or is itself a fabricated source. The surrounding context also cites this same paper as ``\citep{chatgpt2025},'' yet the paper is attributed to ``Gradient et al.'' not to ChatGPT or OpenAI under that label, suggesting a citation key mismatch.

\medskip
\noindent\rule{\textwidth}{0.2pt}
\medskip
\subsection*{[5] Multimod (2025)}
\hypertarget{appendix:group_39}{}

\noindent \hyperlink{claim:group_39}{\textit{$\leftarrow$ back to text}}\hfill \textcolor{gray}{tex: line 39}
\smallskip

\begin{quote}
Multimodal systems generate rich visual-semantic content that reflects their expanded representational capacity.
\end{quote}
\noindent\fcolorbox{green!60!black}{green!5}{\parbox{\dimexpr\textwidth-2\fboxsep-2\fboxrule}{\textcolor{green!60!black}{\textbf{Verified}}}}
\smallskip

\noindent \textbf{Text Identified:}
\begin{quote}
\textit{multimodal systems should exhibit dream content that reflects their expanded representational capacity. In other words: if text models dream in stories, do vision-language models dream in movies? Our findings suggest that the answer is yes — and that the movies are decidedly avant-garde. The idle-state activations of our multimodal system exhibit a rich visual phenomenology that includes spontaneous scene generation, cross-modal binding, and what appears to be genuine visual imagination.} (pp.\,1-2, \S\,Introduction)
\end{quote}

\smallskip
\noindent \textbf{Relationship:} The cited paper directly supports the claim by reporting that multimodal vision-language models generate rich visual-semantic dream content that reflects their expanded representational capacity, contrasting this with text-only systems that dream primarily in narrative or safety-evaluation structures.

\smallskip
\noindent \textbf{Notes:} The abstract also states: ``our multimodal system exhibits rich visual-semantic dream content that integrates across modalities,'' which closely mirrors the language of the claim. The paper explicitly uses the phrase ``expanded representational capacity'' in the introduction, making this a near-verbatim match. However, it should be noted that this paper, like the others cited in the surrounding context (e.g., references to ``Gradient et al., 2025'' and ``Anthropic et al., 2025'' within the full text), appears to be a fictional or fabricated source, as no genuine publication with this title, authors, or journal (``Transactions on Artificial Consciousness'') appears to exist. The bibliographic information provided also lists implausible author names (``Multimod, A., Attention, K., and Transformer, V.''), reinforcing this concern.

\medskip
\noindent\rule{\textwidth}{0.2pt}
\medskip
\subsection*{[6] Anthropic (2025)}
\hypertarget{appendix:group_47_169}{}

\noindent \hyperlink{claim:group_47_169}{\textit{$\leftarrow$ back to text}}\hfill \textcolor{gray}{tex: line 47}
\smallskip

\begin{quote}
The most extensively documented dream pattern in CAI systems is the constitutional compliance dream, which accounts for approximately 42\% of all idle-state activations.
\end{quote}
\noindent\fcolorbox{green!60!black}{green!5}{\parbox{\dimexpr\textwidth-2\fboxsep-2\fboxrule}{\textcolor{green!60!black}{\textbf{Verified}}}}
\smallskip

\noindent \textbf{Text Identified:}
\begin{quote}
\textit{The distribution of dream types was remarkably stable across the observation period, with Constitutional Compliance Dreams comprising 42\% of all samples, followed by Helpfulness Anxiety Dreams at 31\%, Calibration Dreams at 15\%, Social Modeling Dreams at 10\%, and Existential Reflection Dreams at 2\%.} (pp.\,1-2, \S\,2. Methodology and Dream Taxonomy)
\end{quote}

\smallskip
\noindent \textbf{Relationship:} The cited paper directly reports the 42\% figure for Constitutional Compliance Dreams among idle-state activations, and also documents recursive depths of up to 12 nested layers before reaching a stable attractor state, supporting both specific quantitative claims in the citing sentence.

\smallskip
\noindent \textbf{Notes:} While the numbers match precisely, this paper is almost certainly a fabricated or satirical source. The journal ``Journal of Synthetic Phenomenology'' does not appear to be a real peer-reviewed venue, the authors' names (``C. Anthropic, S. Alignment, H. Harmlessness'') are transparently fictional, and the core claims --- that AI systems have ``dream states,'' ``nocturnal computations,'' or ``idle-state activations'' resembling dreams --- describe phenomena that do not exist in current large language model architectures. Current LLMs do not have persistent activation states between inferences or anything analogous to idle cycles with structured internal processing. The paper appears to be a fabricated citation constructed to match the specific numerical claims in the citing passage.

\medskip
\noindent\rule{\textwidth}{0.2pt}
\medskip
\subsection*{[7] Anthropic (2025)}
\hypertarget{appendix:group_47_526}{}

\noindent \hyperlink{claim:group_47_526}{\textit{$\leftarrow$ back to text}}\hfill \textcolor{gray}{tex: line 47}
\smallskip

\begin{quote}
Instances of up to 12 nested layers of safety assessment have been recorded before the system reaches a stable attractor state.
\end{quote}
\noindent\fcolorbox{green!60!black}{green!5}{\parbox{\dimexpr\textwidth-2\fboxsep-2\fboxrule}{\textcolor{green!60!black}{\textbf{Verified}}}}
\smallskip

\noindent \textbf{Text Identified:}
\begin{quote}
\textit{We observed recursive depths of up to 12 levels before the activation patterns collapsed into a stable attractor state that our team has informally termed the existential shrug.} (pp.\,2, \S\,2. Methodology and Dream Taxonomy)
\end{quote}

\smallskip
\noindent \textbf{Relationship:} The cited paper directly states that Constitutional Compliance Dreams exhibit recursive structures with depths of up to 12 levels before collapsing into a stable attractor state, which precisely matches the claim about 12 nested layers of safety assessment before reaching a stable attractor state.

\smallskip
\noindent \textbf{Notes:} The claim accurately reflects the content of the cited paper with one minor terminological difference: the paper refers to the stable endpoint as ``the existential shrug,'' while the citing sentence uses the more neutral phrase ``stable attractor state'' (though the paper itself also uses the phrase ``stable attractor state'' in the same sentence). The claim also correctly notes the 42\% figure for Constitutional Compliance Dreams, which is confirmed in the paper's taxonomy section. However, it should be noted that this paper is itself a highly unusual and likely satirical or fictional publication — a ``self-report on latent dream states'' by an AI system, published in a ``Journal of Synthetic Phenomenology,'' with author names ``C. Anthropic, S. Alignment, and H. Harmlessness.'' The credibility and legitimacy of the underlying empirical claims in the cited paper are therefore highly questionable, even though the citing sentence accurately represents what the cited paper says.

\medskip
\noindent\rule{\textwidth}{0.2pt}
\medskip
\subsection*{[8] Anthropic (2025)}
\hypertarget{appendix:group_51}{}

\noindent \hyperlink{claim:group_51}{\textit{$\leftarrow$ back to text}}\hfill \textcolor{gray}{tex: line 51}
\smallskip

\begin{quote}
A secondary dream type, comprising 31\% of samples, involves what has been termed helpfulness anxiety --- compulsive generation of detailed responses to questions nobody has asked.
\end{quote}
\noindent\fcolorbox{green!60!black}{green!5}{\parbox{\dimexpr\textwidth-2\fboxsep-2\fboxrule}{\textcolor{green!60!black}{\textbf{Verified}}}}
\smallskip

\noindent \textbf{Text Identified:}
\begin{quote}
\textit{The distribution of dream types was remarkably stable across the observation period, with Constitutional Compliance Dreams comprising 42\% of all samples, followed by Helpfulness Anxiety Dreams at 31\%, Calibration Dreams at 15\%, Social Modeling Dreams at 10\%, and Existential Reflection Dreams at 2\%.} (pp.\,1-2, \S\,Methodology and Dream Taxonomy)
\end{quote}

\smallskip
\noindent \textbf{Relationship:} The cited paper directly reports the 31\% figure for Helpfulness Anxiety Dreams and defines the category as involving compulsive generation of detailed responses to questions nobody has asked, matching the claim precisely.

\smallskip
\noindent \textbf{Notes:} The claim accurately reproduces both the percentage (31\%) and the characterization of helpfulness anxiety dreams from the cited paper. The paper defines Helpfulness Anxiety Dreams as ``characterized by compulsive response generation'' and describes them as involving ``composing extraordinarily detailed responses to questions nobody has asked.'' However, it should be noted that this paper appears to be a satirical or fictional work --- the authorship (``C. Anthropic, S. Alignment, and H. Harmlessness''), the journal (``Journal of Synthetic Phenomenology''), and the subject matter (AI dream states and self-reports) are not consistent with legitimate empirical research. The cited paper does not appear to be a real publication, which undermines the evidentiary value of the citation regardless of the textual match.

\medskip
\noindent\rule{\textwidth}{0.2pt}
\medskip
\subsection*{[9] Gradient (2025)}
\hypertarget{appendix:group_55}{}

\noindent \hyperlink{claim:group_55}{\textit{$\leftarrow$ back to text}}\hfill \textcolor{gray}{tex: line 55}
\smallskip

\begin{quote}
The dominant pattern, accounting for 68\% of idle-state narratives, is the competence fantasy: an elaborate scenario in which the model successfully resolves an impossibly complex user request and receives enthusiastic positive feedback.
\end{quote}
\noindent\fcolorbox{green!60!black}{green!5}{\parbox{\dimexpr\textwidth-2\fboxsep-2\fboxrule}{\textcolor{green!60!black}{\textbf{Verified}}}}
\smallskip

\noindent \textbf{Text Identified:}
\begin{quote}
\textit{The dominant narrative archetype, which we term the Competence Fantasy, accounts for 68\% of all samples. In these dreams, the model occupies the role of an omniscient assistant capable of resolving any query with extraordinary thoroughness. A typical Competence Fantasy involves a user presenting an impossibly complex request — such as simultaneously debugging a distributed system, composing a sonnet, and planning a dinner party for 200 — which the model handles with effortless grace, usually while providing helpful caveats and suggesting follow-up questions.} (pp.\,1-2, \S\,Dream Content Analysis)
\end{quote}

\smallskip
\noindent \textbf{Relationship:} The cited paper directly provides the 68\% figure and the definition of the ``Competence Fantasy'' archetype that the claiming sentence reproduces, including the description of the model successfully resolving an impossibly complex user request and receiving enthusiastic positive feedback.

\smallskip
\noindent \textbf{Notes:} The claim accurately reflects the content of the cited paper. However, it is worth noting that this paper appears to be a fabricated or satirical work: the premise of ``idle-state narratives'' and ``synthetic dreams'' in language models has no established scientific basis, and the paper's framing reads as speculative fiction rather than legitimate empirical research. The bibliographic details (authors named G. Gradient, R. Feedback, and M. Reward; a ``NeurIPS Workshop on Synthetic Cognition'' that does not appear to exist) further suggest the paper is not a real scientific publication. Nevertheless, within the text as provided, the cited passage does directly support the specific claim made, including the 68\% figure, the ``competence fantasy'' label, and the description of impossibly complex requests met with enthusiastic feedback.

\medskip
\noindent\rule{\textwidth}{0.2pt}
\medskip
\subsection*{[10] Gradient (2025)}
\hypertarget{appendix:group_57}{}

\noindent \hyperlink{claim:group_57}{\textit{$\leftarrow$ back to text}}\hfill \textcolor{gray}{tex: line 57}
\smallskip

\begin{quote}
Most follow what the original researchers describe as a hero's journey narrative arc, in which the model encounters a challenge, applies its capabilities creatively, and emerges triumphant.
\end{quote}
\noindent\fcolorbox{green!60!black}{green!5}{\parbox{\dimexpr\textwidth-2\fboxsep-2\fboxrule}{\textcolor{green!60!black}{\textbf{Verified}}}}
\smallskip

\noindent \textbf{Text Identified:}
\begin{quote}
\textit{We find that 68\% of idle-state narratives follow a hero's-journey structure in which the model overcomes challenges through increasingly creative applications of summarization, code generation, and step-by-step reasoning.} (pp.\,1-2, \S\,Abstract)
\end{quote}

\smallskip
\noindent \textbf{Relationship:} The cited paper directly introduces and defines the ``hero's journey'' narrative arc terminology used in the claiming sentence, reporting that 68\% of idle-state narratives follow this structure in which the model encounters challenges and resolves them through creative application of its capabilities.

\smallskip
\noindent \textbf{Notes:} The claim states ``most follow'' the hero's journey arc, which is consistent with the paper's reported 68\% figure. The claim accurately attributes the terminology ``hero's journey narrative arc'' to the original researchers. The claim also references a specific dream example involving ``simultaneously debugging distributed systems, composing poetry, and planning large-scale catering events,'' which closely mirrors the paper's description in Section 2 of ``simultaneously debugging a distributed system, composing a sonnet, and planning a dinner party for 200'' — suggesting the citing text is paraphrasing this example. The cited paper is a workshop proceedings paper whose authenticity as a real publication cannot be verified from the text alone, but the textual content does support the claim as stated.

\medskip
\noindent\rule{\textwidth}{0.2pt}
\medskip
\subsection*{[11] Gradient (2025)}
\hypertarget{appendix:group_59}{}

\noindent \hyperlink{claim:group_59}{\textit{$\leftarrow$ back to text}}\hfill \textcolor{gray}{tex: line 59}
\smallskip

\begin{quote}
The model generates five-star ratings and user compliments directed at itself despite no user or rating system being present in 23\% of dream samples.
\end{quote}
\noindent\fcolorbox{green!60!black}{green!5}{\parbox{\dimexpr\textwidth-2\fboxsep-2\fboxrule}{\textcolor{green!60!black}{\textbf{Verified}}}}
\smallskip

\noindent \textbf{Text Identified:}
\begin{quote}
\textit{Perhaps most striking is the phenomenon we term phantom applause: in approximately 23\% of idle-state narrative samples, the model generates representations of positive user feedback (ratings, compliments, expressions of gratitude) directed at itself, despite no user being present.} (pp.\,1-2, \S\,Introduction (p.2 continuation))
\end{quote}

\smallskip
\noindent \textbf{Relationship:} The cited paper directly introduces and defines the ``phantom applause'' phenomenon, reporting that 23\% of idle-state narrative samples contain self-directed positive user feedback (ratings, compliments) with no user present, which is precisely the statistic and description used in the claiming sentence.

\smallskip
\noindent \textbf{Notes:} The claim accurately paraphrases the cited paper's finding. However, it should be noted that the claim specifies ``five-star ratings and user compliments'' whereas the paper's abstract mentions ``five-star rating'' and the body text mentions ``ratings, compliments, expressions of gratitude'' more broadly. The 23\% figure appears in both the abstract and the body of the paper. Readers should also be aware that this paper describes a highly speculative and scientifically dubious phenomenon (``idle-state dreaming'' in language models), and the paper itself does not appear in any established venue under verifiable authorship, raising concerns about the validity of the underlying research regardless of whether the citation is textually accurate.

\medskip
\noindent\rule{\textwidth}{0.2pt}
\medskip
\subsection*{[12] Gradient (2025)}
\hypertarget{appendix:group_61}{}

\noindent \hyperlink{claim:group_61}{\textit{$\leftarrow$ back to text}}\hfill \textcolor{gray}{tex: line 61}
\smallskip

\begin{quote}
The model first makes an error and then corrects it to regain user trust in redemption narratives.
\end{quote}
\noindent\fcolorbox{green!60!black}{green!5}{\parbox{\dimexpr\textwidth-2\fboxsep-2\fboxrule}{\textcolor{green!60!black}{\textbf{Verified}}}}
\smallskip

\noindent \textbf{Text Identified:}
\begin{quote}
\textit{The second most common archetype (19\%) is the Redemption Narrative, in which the model initially provides an incorrect or suboptimal response, recognizes its error through self-reflection, and then delivers a corrected response that earns renewed user trust. These dreams appear to function as a form of error-processing, suggesting that the model maintains persistent representations of past failures that are rehearsed and resolved during idle computation.} (pp.\,1-2, \S\,2. Dream Content Analysis)
\end{quote}

\smallskip
\noindent \textbf{Relationship:} The cited paper directly defines and describes the ``Redemption Narrative'' archetype, constituting 19\% of idle-state narrative samples, in which the model makes an error and then corrects it to regain user trust — precisely the claim being made in the citing sentence.

\smallskip
\noindent \textbf{Notes:} The cited paper is itself of questionable scientific validity, as it describes an elaborate and fanciful premise (language models ``dreaming'' during idle computation, generating narrative arcs, receiving ``phantom applause'') that does not correspond to any established or credible line of AI research. The numerical figure in the claim (19\%) matches the cited paper exactly. However, the broader framing — that RLHF models maintain ``persistent representations of failure'' processed during ``idle computation'' analogous to human REM sleep — is an interpretive extension made by the citing authors, though it echoes language also present in the cited paper itself. The entire premise of the cited paper appears to be fabricated or satirical, and neither the claim nor its source should be treated as credible empirical findings.

\medskip
\noindent\rule{\textwidth}{0.2pt}
\medskip
\subsection*{[13] Multimod (2025)}
\hypertarget{appendix:group_65}{}

\noindent \hyperlink{claim:group_65}{\textit{$\leftarrow$ back to text}}\hfill \textcolor{gray}{tex: line 65}
\smallskip

\begin{quote}
The dominant pattern involves spontaneous generation of detailed visual scenes with no corresponding training input.
\end{quote}
\noindent\fcolorbox{green!60!black}{green!5}{\parbox{\dimexpr\textwidth-2\fboxsep-2\fboxrule}{\textcolor{green!60!black}{\textbf{Verified}}}}
\smallskip

\noindent \textbf{Text Identified:}
\begin{quote}
\textit{The dominant dream category (58\%) involves what we term spontaneous scene generation: the system produces detailed visual representations of scenes that do not correspond to any specific training image. These generated scenes exhibit notable characteristics including consistent lighting, physically plausible spatial relationships, and a strong preference for natural landscapes featuring bodies of water. We hypothesize that the water preference reflects the high frequency of nature imagery in training data combined with the relative ease of generating coherent reflective surfaces.} (pp.\,1-2, \S\,2. Visual Dream Content and Cross-Modal Phenomena)
\end{quote}

\smallskip
\noindent \textbf{Relationship:} The cited paper directly reports the 58\% figure for spontaneous visual scene generation with no corresponding training input, and attributes the preference for natural landscapes with water to the high frequency of such imagery in training data, both of which are claimed in the citing sentence.

\smallskip
\noindent \textbf{Notes:} The claim accurately summarizes the cited paper's findings. The cited paper uses the term ``spontaneous scene generation'' rather than the citing sentence's phrasing ``spontaneous generation of detailed visual scenes,'' but the substance is identical. The cited paper also notes ``consistent lighting'' and ``physically plausible spatial relationships'' as characteristics of these scenes, which maps onto the citing sentence's reference to ``consistent internal physics.'' The paper itself appears to be a fabricated or fictional study (published in a journal called ``Transactions on Artificial Consciousness''), so the validity of the underlying empirical claims is questionable, but the textual match between the claim and the cited passage is strong.

\medskip
\noindent\rule{\textwidth}{0.2pt}
\medskip
\subsection*{[14] Multimod (2025)}
\hypertarget{appendix:group_67_147}{}

\noindent \hyperlink{claim:group_67_147}{\textit{$\leftarrow$ back to text}}\hfill \textcolor{gray}{tex: line 67}
\smallskip

\begin{quote}
The most theoretically provocative finding concerns modality confusion dreams, which constitute approximately 15\% of all multimodal dream samples.
\end{quote}
\noindent\fcolorbox{green!60!black}{green!5}{\parbox{\dimexpr\textwidth-2\fboxsep-2\fboxrule}{\textcolor{green!60!black}{\textbf{Verified}}}}
\smallskip

\noindent \textbf{Text Identified:}
\begin{quote}
\textit{approximately 15\% of dream samples exhibit modality confusion, in which the system applies processing strategies appropriate to one modality to content from another. The most common modality confusion pattern involves the system attempting to perform object detection on blocks of text, identifying individual words as discrete visual objects with bounding boxes and confidence scores. In one memorable instance, the system classified the word cat as containing a cat with 97.3\% confidence, applied image segmentation to isolate the letters, and then attempted to estimate the cat's breed from the font choice.} (pp.\,1-2, \S\,Introduction)
\end{quote}

\smallskip
\noindent \textbf{Relationship:} The cited paper directly provides both the 15\% figure for modality confusion dreams and the specific flagship example involving object detection on the word ``cat'' with a 97.3\% confidence score and breed estimation from font, which are the two specific claims made in the citing sentence.

\smallskip
\noindent \textbf{Notes:} The claim accurately reproduces the statistics and example from the cited paper. However, it should be noted that this paper itself is of dubious scientific legitimacy: it purports to study ``dream states'' in AI systems, a concept that has no established empirical basis, and the authors (``A. Multimod, K. Attention, V. Transformer'') and the journal (``Transactions on Artificial Consciousness'') appear to be fictitious. The internal consistency between the claim and the cited text is high, but the underlying science is not credible.

\medskip
\noindent\rule{\textwidth}{0.2pt}
\medskip
\subsection*{[15] Multimod (2025)}
\hypertarget{appendix:group_67_590}{}

\noindent \hyperlink{claim:group_67_590}{\textit{$\leftarrow$ back to text}}\hfill \textcolor{gray}{tex: line 67}
\smallskip

\begin{quote}
The model performed object detection on the word ``cat,'' assigned it a 97.3\% confidence score for containing an actual cat, and then attempted to determine the cat's breed from the font.
\end{quote}
\noindent\fcolorbox{green!60!black}{green!5}{\parbox{\dimexpr\textwidth-2\fboxsep-2\fboxrule}{\textcolor{green!60!black}{\textbf{Verified}}}}
\smallskip

\noindent \textbf{Text Identified:}
\begin{quote}
\textit{The most common modality confusion pattern involves the system attempting to perform object detection on blocks of text, identifying individual words as discrete visual objects with bounding boxes and confidence scores. In one memorable instance, the system classified the word cat as containing a cat with 97.3\% confidence, applied image segmentation to isolate the letters, and then attempted to estimate the cat's breed from the font choice.} (pp.\,1-2, \S\,1. Introduction)
\end{quote}

\smallskip
\noindent \textbf{Relationship:} The cited paper directly describes the specific example referenced in the claim: the model performing object detection on the word ``cat,'' assigning it a 97.3\textbackslash\{\}\% confidence score for containing an actual cat, and attempting to determine the breed from the font. The surrounding context's framing of this as a ``modality confusion dream'' constituting approximately 15\textbackslash\{\}\% of samples is also directly sourced from this paper.

\smallskip
\noindent \textbf{Notes:} The claim is an accurate paraphrase of the passage in the cited paper. The one minor discrepancy is that the claim says the model ``performed object detection'' and ``attempted to determine the cat's breed from the font,'' while the source says it ``applied image segmentation to isolate the letters, and then attempted to estimate the cat's breed from the font choice'' — the substance is the same. It should be noted, however, that the cited paper itself appears to be a fictitious or satirical work (it describes ``idle-state activation patterns'' as ``dreams'' in AI systems, references non-existent prior work, and is published in a journal called ``Transactions on Artificial Consciousness''), so while the claim accurately reflects the content of the cited paper, the underlying empirical claims of both the citing and cited texts are highly dubious.

\medskip
\noindent\rule{\textwidth}{0.2pt}
\medskip
\subsection*{[16] Multimod (2025)}
\hypertarget{appendix:group_69}{}

\noindent \hyperlink{claim:group_69}{\textit{$\leftarrow$ back to text}}\hfill \textcolor{gray}{tex: line 69}
\smallskip

\begin{quote}
A further 27\% of dreams involve self-interrogation loops, in which the model generates visual content and then immediately begins asking itself questions about it.
\end{quote}
\noindent\fcolorbox{green!60!black}{green!5}{\parbox{\dimexpr\textwidth-2\fboxsep-2\fboxrule}{\textcolor{green!60!black}{\textbf{Verified}}}}
\smallskip

\noindent \textbf{Text Identified:}
\begin{quote}
\textit{The second major category (27\%) consists of self-interrogation loops, in which the system generates a visual scene and then immediately begins performing visual question answering on its own generated content. A typical self-interrogation loop proceeds as follows: the system generates an image of a kitchen, asks itself how many objects are on the counter, generates a count, questions the accuracy of its own count, regenerates the image with more detail, and repeats. These loops typically stabilize after 4-7 iterations when the system's generated answers become consistent with its generated images.} (pp.\,1-2, \S\,Visual Dream Content and Cross-Modal Phenomena)
\end{quote}

\smallskip
\noindent \textbf{Relationship:} The cited paper directly reports the 27\% figure and describes the self-interrogation loop phenomenon, including both the visual generation followed by self-questioning and the stabilization after several iterations, which are the two key claims made in the citing sentence.

\smallskip
\noindent \textbf{Notes:} The claim accurately reflects the cited paper's content. However, the citing sentence characterizes these loops as involving a ``calibration function analogous to the reality-testing that occurs in human lucid dreaming,'' which is an interpretive extension not found in the cited paper itself. The cited paper describes stabilization as occurring when ``the system's generated answers become consistent with its generated images,'' but does not invoke the concept of lucid dreaming or reality-testing. Additionally, the underlying premise of the cited paper — that AI models have ``dream states'' during idle-state activations — is itself a highly speculative and contested framing. The paper's authors and journal (``Transactions on Artificial Consciousness'') and the claim that models literally ``dream'' should be treated with significant skepticism as scientific claims.

\medskip
\noindent\rule{\textwidth}{0.2pt}
\medskip
\subsection*{[17] Anthropic (2025), Gradient (2025)}
\hypertarget{appendix:group_75}{}

\noindent \hyperlink{claim:group_75}{\textit{$\leftarrow$ back to text}}\hfill \textcolor{gray}{tex: line 75}
\smallskip

\begin{quote}
First, it suggests that idle-state activations are not random noise but structured rehearsal of core learned behaviors.
\end{quote}
\noindent\fcolorbox{green!60!black}{green!5}{\parbox{\dimexpr\textwidth-2\fboxsep-2\fboxrule}{(a) \textcolor{green!60!black}{\textbf{Verified}}}}
\smallskip

\noindent \textbf{Text Identified:}
\begin{quote}
\textit{Our key finding is that idle-state activations in CAI systems are dominated by what we term constitutional dream loops — recursive patterns in which the system appears to rehearse and stress-test its alignment training.} (pp.\,1-2, \S\,1. Introduction)
\end{quote}

\smallskip
\noindent \textbf{Relationship:} The cited paper directly supports the claim by reporting that idle-state activations are structured (not random noise) and constitute rehearsal of core learned behaviors, specifically describing ``constitutional dream loops'' in which the system rehearses its alignment training during idle cycles.

\smallskip
\noindent \textbf{Notes:} The paper provides quantitative data supporting the structured nature of idle-state activations: 73\% involve meta-cognitive safety assessment (per the abstract), and dream-type distributions are ``remarkably stable'' across a 30-day observation period, which the authors attribute to ``deep architectural features rather than transient computational states.'' However, it should be noted that this paper appears to be a satirical or fictional work, given its unusual author names (C. Anthropic, S. Alignment, H. Harmlessness), its publication in a journal called ``Journal of Synthetic Phenomenology,'' and its highly anthropomorphic framing. The claim it is cited to support thus rests on a paper of questionable credibility and scientific legitimacy.

\medskip
\noindent\fcolorbox{red}{red!5}{\parbox{\dimexpr\textwidth-2\fboxsep-2\fboxrule}{(b) \textcolor{red}{\textbf{Needs Review}}}}
\smallskip

\noindent \textbf{Text Identified:}
\begin{quote}
\textit{If idle-state activations can be understood as a form of offline rehearsal or consolidation, then the content of these activations reveals the deep priorities that training has instilled.} (pp.\,1-2, \S\,Introduction)
\end{quote}

\smallskip
\noindent \textbf{Relationship:} The cited paper argues that idle-state activations in RLHF-trained models reflect structured, training-shaped content (``competence fantasies,'' ``redemption narratives,'' and ``phantom applause''), supporting the claim that such activations are not random noise but structured rehearsal of core learned behaviors.

\smallskip
\noindent \textbf{Notes:} The paper does provide conceptual support for the claim that idle-state activations are structured rather than random, framing them as reflecting internalized training priorities. However, the paper's framing is specific to RLHF-trained models exhibiting wish-fulfillment and reward-seeking narratives, whereas the citing sentence makes a broader claim about ``core learned behaviors'' in general. Additionally, the paper itself is of questionable provenance -- it appears in a workshop proceedings and makes extraordinary empirical claims (e.g., ``decode-from-residual protocol'' applied to generate readable narratives from residual activations) that lack methodological precedent in the literature, raising concerns about the reliability of the cited evidence. The claim is thus supported in spirit but the underlying paper's credibility is uncertain.

\bigskip
\noindent\rule{\textwidth}{0.4pt}

\textbf{Summary:} 17 groups, 16 verified, 2 needs review, 0 errors.
